Self-improving systems---reflection loops, memory compilers, invariant
discovery, and self-refining harnesses---may generate candidates for
change. The boundary we enforce is that none of these paths can produce
authority.

\begin{theorem}[Candidate-only cannot authorize]\label{thm:candidate-only}
A candidate artifact cannot, by itself, produce an authority grant.
\end{theorem}

\begin{proof}
By construction of the authority-grant relation: a grant is admissible only
when its root identifier differs from the candidate's emission path. The
Lean statement \texttt{candidate\_only\_cannot\_authorize} makes this a
type-level obligation, and the reference gate
(\texttt{promote\_candidate}) returns REJECT when the candidate is its own
root.
\end{proof}

\begin{theorem}[No same-epoch self-certification]\label{thm:epoch}
Reflection generated in epoch $E$ cannot authorize an action in epoch $E$.
\end{theorem}

\begin{proof}
Epoch separation is enforced at the reflection boundary: the reflection
epoch must strictly precede the authority-eligibility epoch
($E$ execution $\to$ $E{+}1$ reflection $\to$ $E{+}2$ authority). The Lean
statement \texttt{no\_same\_epoch\_self\_certification} and the executable
gate \texttt{same\_epoch\_self\_certification} both reject
$\texttt{reflectionEpoch} = \texttt{authorityEpoch}$.
\end{proof}

The estate's reflection module exposes \texttt{assert\_reflection\_only},
which raises on actions in the \texttt{authorize}/\texttt{grant}/
\texttt{promote}/\texttt{admit}/\texttt{execute} classes. Reflection paths
may observe, score, challenge, and recommend; they may not act. This is a
runtime guard in addition to the type-level obligation.
